\documentclass[11pt,a4paper]{article}
\usepackage[utf8]{inputenc}
\usepackage{amsmath,amssymb,amsthm}
\usepackage{mathtools}
\usepackage{bbm}
\usepackage{booktabs}
\usepackage[dvipsnames]{xcolor}
\usepackage{tcolorbox}
\usepackage[margin=2.5cm]{geometry}
\usepackage{hyperref}
\usepackage{enumitem}
\usepackage{cleveref}

\tcbuselibrary{skins,breakable}

% Theorem environments
\newtheorem{theorem}{Theorem}
\newtheorem{proposition}[theorem]{Proposition}
\newtheorem{lemma}[theorem]{Lemma}
\newtheorem{corollary}[theorem]{Corollary}
\newtheorem{definition}[theorem]{Definition}
\newtheorem{assumption}{Assumption}
\newtheorem{remark}[theorem]{Remark}

% Notation shortcuts
\newcommand{\R}{\mathbb{R}}
\newcommand{\E}{\mathbb{E}}
\newcommand{\N}{\mathbb{N}}
\newcommand{\Gr}{\mathrm{Gr}}
\newcommand{\tr}{\mathrm{tr}}
\newcommand{\diag}{\mathrm{diag}}
\newcommand{\reg}{\mathrm{Reg}}
\newcommand{\OPT}{\mathrm{OPT}}
\newcommand{\col}{\mathrm{col}}
\newcommand{\sign}{\mathrm{sign}}
\newcommand{\norm}[1]{\left\|#1\right\|}
\newcommand{\inner}[2]{\langle #1, #2 \rangle}
\newcommand{\dG}{d_{\Gr}}
\DeclareMathOperator*{\argmin}{arg\,min}
\DeclareMathOperator*{\argmax}{arg\,max}
\DeclareMathOperator{\Span}{span}
\DeclareMathOperator{\rank}{rank}

\title{Formal Proofs: Dual-Timescale Adaptation for\\Streaming Low-Rank Learning under Multi-Rate Drift}
\author{Supplementary Material}
\date{April 2026}

\begin{document}
\maketitle

\begin{abstract}
This supplement contains complete formal proofs for all theoretical claims
in the paper \emph{Dual-Timescale Adaptation Is Necessary and Sufficient
for Streaming Low-Rank Learning under Multi-Rate Drift}.
We provide: (i)~a precise problem setting with explicit assumptions
(\Cref{sec:setting}), (ii)~a necessity lower bound for monotone-step
learners (\Cref{sec:theorem1}), (iii)~an upper bound for the dual-timescale
algorithm under both general variation and jump-process regimes
(\Cref{sec:theorem2}), and (iv)~three corollaries on optimal merging,
consolidation frequency, and selective consolidation (\Cref{sec:corollaries}).
\end{abstract}

\tableofcontents
\newpage

%% setting.tex — Problem setting, assumptions, definitions
%% Fixes Issues: #1 (no proofs), #2 (strong convexity), #3/#4 (gauge),
%%               #12 (regret definition), #13 (probability mode)
\section{Problem Setting and Assumptions}
\label{sec:setting}

%% ======================================================================
\subsection{Streaming Low-Rank Learning}
%% ======================================================================

We consider a learner receiving a stream of data $(x_t, y_t)_{t=1}^T$
drawn i.i.d.\ from a time-varying distribution $\mathcal{D}_t$. The
learner maintains a parameter $\theta_t \in \R^d$ and suffers per-round
loss $\ell_t(\theta) \coloneqq \ell(\theta; x_t, y_t)$.

\begin{definition}[Population loss]
\label{def:population-loss}
The \emph{population loss} at time $t$ is
\[
  L_t(\theta) \coloneqq \E_{(x,y) \sim \mathcal{D}_t}\bigl[\ell(\theta; x, y)\bigr].
\]
\end{definition}

\begin{definition}[Ground-truth optimum]
\label{def:optimum}
The time-varying optimum is $\theta^*_t \coloneqq \argmin_{\theta \in \R^d} L_t(\theta)$,
assumed to exist and be unique (ensured by strong convexity below).
\end{definition}

%% ======================================================================
\subsection{Two-Rate Structural Decomposition}
%% ======================================================================

\begin{assumption}[Low-rank decomposition]
\label{ass:decomposition}
There exist a slowly varying subspace $\mathcal{U}_t = \col(U_t)$ where
$U_t \in \R^{d \times r}$ has orthonormal columns ($U_t^\top U_t = I_r$,
so $U_t$ represents a point on the Grassmannian $\Gr(d,r)$), a fast-varying
coefficient vector $a_t \in \R^r$, and a residual $\rho_t \in \R^d$
such that
\[
  \theta^*_t = U_t\, a_t + \rho_t,
  \qquad
  \norm{\rho_t} \leq \rho_{\max},
  \qquad
  1 \leq r < d.
\]
The residual $\rho_t$ captures the approximation error of projecting
$\theta^*_t$ onto the rank-$r$ subspace.
\end{assumption}

%% ======================================================================
\subsection{Gauge-Invariant Subspace Distance}
%% ======================================================================

\begin{remark}[Gauge invariance --- fixes Issues \#3, \#4]
\label{rem:gauge}
The Frobenius distance $\norm{U_t - U_{t'}}_F$ is \emph{not}
gauge-invariant: for any orthogonal $R \in \mathcal{O}(r)$,
$U_t' = U_t R$ represents the \emph{same} subspace but can have
$\norm{U_t - U_t'}_F$ up to $\sqrt{2r}$.
We therefore use the \emph{principal-angle (geodesic) distance}
on $\Gr(d,r)$.
\end{remark}

\begin{definition}[Grassmann distance]
\label{def:grassmann-dist}
Let $U, U' \in \R^{d \times r}$ have orthonormal columns. Let
$\sigma_1 \geq \cdots \geq \sigma_r$ be the singular values of
$U^\top U'$. The \emph{principal angles} are
$\theta_i = \arccos(\min(\sigma_i, 1))$ for $i = 1, \ldots, r$.
The Grassmann distance is
\[
  \dG(\col(U),\, \col(U'))
  \coloneqq \norm{\sin \Theta}_F
  = \Bigl(\sum_{i=1}^r \sin^2\theta_i\Bigr)^{1/2}.
\]
Equivalently, $\dG^2 = r - \norm{U^\top U'}_F^2$ when both have
orthonormal columns.
\end{definition}

\begin{lemma}[Grassmann distance properties]
\label{lem:grassmann-props}
$\dG$ is a metric on $\Gr(d,r)$ satisfying:
\begin{enumerate}[nosep]
  \item $0 \leq \dG \leq \sqrt{r}$;
  \item $\dG(\col(U), \col(UR)) = 0$ for all $R \in \mathcal{O}(r)$
        \textup{(gauge invariance)};
  \item $\dG \leq \norm{U - U'}_F \leq \sqrt{2}\, \dG$ when both
        $U, U'$ are orthonormal \textup{(Frobenius sandwich)}.
\end{enumerate}
\end{lemma}

\begin{proof}
(1) follows from $\sin^2\theta_i \in [0,1]$.
(2): $U$ and $UR$ have the same column space, so
$\sigma_i(U^\top(UR)) = \sigma_i(R) = 1$ for all $i$, giving
$\theta_i = 0$.
(3): Using $\norm{U - U'}_F^2 = 2r - 2\tr(U^\top U')$ and the
relation between trace and singular values:
$\tr(U^\top U') \leq \sum_i \sigma_i = \sum_i \cos\theta_i$,
so $\norm{U-U'}_F^2 = 2\sum_i(1-\cos\theta_i) \leq
2\sum_i \sin^2\theta_i / (1+\cos\theta_i) \cdot 2 \leq
2\sum_i \sin^2\theta_i = 2\dG^2$ using $1-\cos\theta \leq
\sin^2\theta$ for $\theta \in [0,\pi/2]$.
The lower bound: $\norm{U-U'}_F^2 \geq 2\sum_i(1-\cos\theta_i)
\geq \sum_i \sin^2\theta_i = \dG^2$ using $1-\cos\theta \geq
\sin^2\theta/2$.
\end{proof}

%% ======================================================================
\subsection{Drift Variation Budgets}
%% ======================================================================

\begin{definition}[Slow variation]
\label{def:slow-variation}
The \emph{slow (subspace) variation} over horizon $T$ is
\[
  V_s(T) \coloneqq \sum_{t=1}^{T-1} \dG\bigl(\col(U_t),\, \col(U_{t+1})\bigr).
\]
\end{definition}

\begin{definition}[Fast variation --- jump process]
\label{def:fast-variation}
The fast (coefficient) variation is modeled as a \emph{piecewise-constant
jump process}: there exist times $0 = \tau_0 < \tau_1 < \cdots < \tau_{K_T} \leq T$
such that $a_t = a_{\tau_k}$ for $t \in [\tau_k, \tau_{k+1})$. We write
$K_T$ for the number of jumps and define
\[
  V_f(T) \coloneqq \sum_{k=1}^{K_T} \norm{a_{\tau_k} - a_{\tau_{k-1}}}.
\]
\end{definition}

\begin{assumption}[Bounded variation]
\label{ass:bounded-variation}
$V_s(T)$ and $V_f(T)$ are finite for all $T$. We also assume
minimum jump size $\Delta_0 > 0$:
\[
  \norm{a_{\tau_k} - a_{\tau_{k-1}}} \geq \Delta_0
  \quad \text{for all } k = 1, \ldots, K_T.
\]
\end{assumption}

\begin{assumption}[Bounded coefficients]
\label{ass:bounded-coeff}
$\norm{a_t} \leq A_{\max}$ for all $t$. This ensures the jump magnitudes
are bounded: $\norm{a_{\tau_k} - a_{\tau_{k-1}}} \leq 2A_{\max}$.
\end{assumption}

%% ======================================================================
\subsection{Loss Assumptions (on Population Loss)}
%% ======================================================================

The following assumptions apply to the \emph{population loss} $L_t$,
\textbf{not} the per-sample loss $\ell_t$. This is a critical distinction
(cf.\ Issue~\#2): the per-sample squared loss
$\ell_t(\theta) = \frac{1}{2}(\theta^\top x_t - y_t)^2$ has rank-1
Hessian $x_t x_t^\top$ for $d > 1$, which is \emph{not} strongly convex.
But the population loss $L_t(\theta) = \frac{1}{2}\E[(\theta^\top x - y)^2]$
has Hessian $\Sigma_t = \E[x_t x_t^\top]$, which is $\mu$-strongly convex
when $\lambda_{\min}(\Sigma_t) \geq \mu > 0$.

\begin{assumption}[Strong convexity of population loss]
\label{ass:strong-convexity}
For all $t \in [T]$, the population loss $L_t$ is $\mu$-strongly convex:
\[
  L_t(\theta') \geq L_t(\theta) + \inner{\nabla L_t(\theta)}{\theta'-\theta}
  + \frac{\mu}{2}\norm{\theta'-\theta}^2,
  \qquad \forall\, \theta, \theta' \in \R^d,
\]
where $\mu = \inf_{t \in [T]} \lambda_{\min}\bigl(\nabla^2 L_t(\theta^*_t)\bigr) > 0$.
\end{assumption}

\begin{assumption}[$L$-smoothness of population loss]
\label{ass:smoothness}
For all $t \in [T]$, $L_t$ is $L$-smooth:
\[
  \norm{\nabla L_t(\theta) - \nabla L_t(\theta')} \leq L\norm{\theta - \theta'},
  \qquad \forall\, \theta, \theta' \in \R^d.
\]
The condition number is $\kappa \coloneqq L/\mu \geq 1$.
\end{assumption}

\begin{assumption}[Bounded noise]
\label{ass:noise}
The stochastic gradient $g_t(\theta) = \nabla \ell_t(\theta)$ is an
unbiased estimator of $\nabla L_t(\theta)$ with bounded variance:
\[
  \E\bigl[\norm{g_t(\theta) - \nabla L_t(\theta)}^2 \bigm| \theta\bigr]
  \leq \sigma^2, \qquad \forall\, \theta \in \R^d.
\]
\end{assumption}

%% ======================================================================
\subsection{Regret Definition}
%% ======================================================================

\begin{definition}[Expected dynamic regret --- fixes Issue~\#12]
\label{def:regret}
The \emph{expected dynamic regret} of an algorithm producing
$\theta_1, \ldots, \theta_T$ (where $\theta_t$ is $\mathcal{F}_{t-1}$-measurable,
$\mathcal{F}_t = \sigma(x_1,y_1,\ldots,x_t,y_t)$) is
\[
  \reg_T \coloneqq \E\Biggl[\sum_{t=1}^T
    \bigl(L_t(\theta_t) - L_t(\theta^*_t)\bigr)\Biggr],
\]
where the expectation is over the randomness in $(x_t, y_t)_{t=1}^T$
and any internal randomness of the algorithm.
\end{definition}

\begin{remark}[Probability mode --- fixes Issue~\#13]
\label{rem:prob-mode}
All bounds in this supplement are stated \emph{in expectation}.
By Markov's inequality, an expected regret bound of $\E[\reg_T] \leq R$
implies a high-probability bound $\Pr[\reg_T > R/\delta] \leq \delta$.
Sharper high-probability bounds (e.g., via Azuma--Hoeffding for bounded
differences) are possible but beyond the scope of these proofs.
\end{remark}

%% ======================================================================
\subsection{Applicability to LLM Fine-Tuning}
%% ======================================================================

\begin{remark}[Theory--practice scope]
\label{rem:scope}
\Cref{ass:strong-convexity,ass:smoothness} hold for quadratic / linear
regression with well-conditioned covariance. They do \emph{not} hold
for LLM next-token prediction loss, which is non-convex in the full
parameter space and only approximately locally convex in the LoRA subspace
near a pre-trained minimum.

The theorems below therefore provide \emph{motivating theory} for the
dual-timescale principle. The Mirror-LoRA algorithm applied to LLMs is
a \emph{theory-inspired heuristic}, not a formally justified instantiation
of the theoretical framework. The empirical evaluation (Section~5 of
the main paper) demonstrates that the dual-timescale principle yields
strong performance on real LLMs despite the gap between theory and practice.
\end{remark}

%% ======================================================================
\subsection{Assumption Summary}
%% ======================================================================

\begin{center}
\begin{tabular}{clcc}
\toprule
\textbf{ID} & \textbf{Assumption} & \textbf{Used by} & \textbf{Status} \\
\midrule
A1 & $L_t$ is $\mu$-strongly convex (\Cref{ass:strong-convexity})
   & Thm~1,2 & Holds for linear-Gaussian \\
A2 & $L_t$ is $L$-smooth (\Cref{ass:smoothness})
   & Thm~1,2 & Holds for linear-Gaussian \\
A3 & Low-rank decomposition (\Cref{ass:decomposition})
   & Thm~1,2 & By construction (synthetic) \\
A4 & Bounded variation (\Cref{ass:bounded-variation})
   & Thm~1,2 & Finite horizon \\
A5 & Bounded coefficients (\Cref{ass:bounded-coeff})
   & Thm~1 & Replacement jumps \\
A6 & Bounded noise (\Cref{ass:noise})
   & Thm~2 & Bounded loss \\
\bottomrule
\end{tabular}
\end{center}

%% theorem1.tex — Necessity: Lower bound for monotone-step learners
%% Fixes Issues: #6 (quantifier), #7 (recovery lemma), #8 (non-overlap), #9 (Delta->0)
\section{Theorem 1: Necessity of Dual-Timescale Adaptation}
\label{sec:theorem1}

%% ======================================================================
\subsection{Monotone-Step Learner Class}
%% ======================================================================

\begin{definition}[Monotone-step learner --- fixes Issue~\#6]
\label{def:monotone-step}
A \emph{monotone-step learner} $\mathcal{A}$ maintains a single iterate
$\theta_t \in \R^d$ updated as
\[
  \theta_{t+1} = \theta_t - \eta_t\, G_t(\theta_t, x_t, y_t),
\]
where:
\begin{enumerate}[nosep]
  \item $G_t$ is a function satisfying
        $\E[G_t(\theta, x_t, y_t)] = \nabla L_t(\theta) + b_t(\theta)$
        with $\norm{b_t(\theta)} \leq \beta$ (bounded bias);
  \item The step size $\eta_t = c \cdot t^{-\alpha}$ for constants
        $c > 0$ and $\alpha \in [1/2, 1]$, and is \emph{non-increasing}:
        $\eta_{t+1} \leq \eta_t$.
\end{enumerate}
This class includes:
\begin{itemize}[nosep]
  \item SGD with decaying step size ($G_t = g_t$, $\beta = 0$, $\alpha \in [1/2,1]$)
  \item Online gradient descent (OGD) with $\eta_t = c/\sqrt{t}$
  \item Mirror descent with decaying step size
\end{itemize}
This class \textbf{excludes}:
\begin{itemize}[nosep]
  \item Adam / AdaGrad (adaptive preconditioner, effectively non-monotone
        per-coordinate step size)
  \item Restarted algorithms / change-point detectors
  \item Expert mixture methods / strongly adaptive algorithms
  \item Constant step-size SGD with periodic averaging
\end{itemize}
\end{definition}

%% ======================================================================
\subsection{Stationary Regret of Monotone-Step Learners}
%% ======================================================================

\begin{lemma}[Stationary regret --- standard]
\label{lem:stationary-regret}
Under \Cref{ass:strong-convexity,ass:smoothness,ass:noise}, a
monotone-step learner with $\alpha = 1/2$ achieves stationary regret
\[
  \E\Biggl[\sum_{t=1}^T \bigl(L_t(\theta_t) - L_t(\theta^*)\bigr)\Biggr]
  = O\!\left(\frac{L\sigma^2}{\mu^2}\sqrt{T} + \frac{L\norm{\theta_1 - \theta^*}^2}{\mu}\right)
\]
on a stationary stream ($\theta^*_t = \theta^*$ for all $t$, $V_s = V_f = 0$).
\end{lemma}

\begin{proof}
This is the standard analysis of SGD on strongly convex functions.
By the descent lemma for $\mu$-strongly convex $L$:
\[
  \E[\norm{\theta_{t+1} - \theta^*}^2]
  \leq (1 - 2\mu\eta_t + L^2\eta_t^2)\E[\norm{\theta_t - \theta^*}^2]
    + \eta_t^2 \sigma^2.
\]
With $\eta_t = c/\sqrt{t}$ and $c \leq 1/(2L)$:
\[
  \E[\norm{\theta_t - \theta^*}^2]
  \leq \frac{C_{\mathrm{stat}}}{t}, \qquad
  C_{\mathrm{stat}} = \max\!\left(\norm{\theta_1-\theta^*}^2,\,
  \frac{c^2\sigma^2}{2\mu c - 1}\right).
\]
By strong convexity, $L(\theta_t) - L(\theta^*) \leq
\frac{L}{2}\norm{\theta_t-\theta^*}^2$, so
\[
  \reg_T^{\mathrm{stat}} \leq \frac{L\, C_{\mathrm{stat}}}{2}
  \sum_{t=1}^T \frac{1}{t}
  = O(\sqrt{T}).
  \qedhere
\]
\end{proof}

%% ======================================================================
\subsection{Recovery Lemma}
%% ======================================================================

\begin{lemma}[Recovery cost per jump --- fixes Issue~\#7]
\label{lem:recovery-cost}
Under \Cref{ass:strong-convexity,ass:smoothness,ass:noise}, suppose
the optimum jumps from $\theta^*$ to $\theta^{*\prime}$ at time $\tau$
with $\norm{\theta^* - \theta^{*\prime}} = \Delta \geq \Delta_0 > 0$.
A monotone-step learner with step size $\eta_t = c \cdot t^{-\alpha}$
($\alpha \in [1/2,1]$) that had converged to
$\E[\norm{\theta_\tau - \theta^*}^2] \leq \epsilon$ satisfies:
\[
  \sum_{s=\tau}^{\tau + T_{\mathrm{rec}}-1}
  \E\bigl[L_s(\theta_s) - L_s(\theta^{*\prime})\bigr]
  \geq \frac{\mu\Delta^2}{4}\cdot
  \frac{c^{-1}\tau^{\alpha}}{1 + \mu c\, \tau^{-\alpha}}
  \geq \frac{\Delta^2}{8c}\, \tau^{\alpha},
\]
where $T_{\mathrm{rec}}$ is the number of steps until
$\E[\norm{\theta_t - \theta^{*\prime}}^2] \leq
\max(\epsilon, c^2\sigma^2\tau^{-\alpha}/\mu)$.
\end{lemma}

\begin{proof}
After the jump, $\E[\norm{\theta_\tau - \theta^{*\prime}}^2]
\geq (\Delta - \sqrt{\epsilon})^2 \geq \Delta^2/4$ (since
$\sqrt{\epsilon} \leq \Delta/2$ for sufficiently converged states).

The one-step contraction with step size $\eta_\tau = c\tau^{-\alpha}$ gives:
\[
  \E[\norm{\theta_{s+1} - \theta^{*\prime}}^2]
  \leq (1 - \mu\eta_s)\E[\norm{\theta_s - \theta^{*\prime}}^2]
  + \eta_s^2\sigma^2.
\]
For $s \in [\tau, \tau + T_{\mathrm{rec}})$, since $\eta_s$ is
non-increasing and $\eta_s \leq \eta_\tau = c\tau^{-\alpha}$, the
contraction rate per step is at most $\mu\eta_\tau = \mu c \tau^{-\alpha}$.

To reduce $\E[\norm{\theta_s - \theta^{*\prime}}^2]$ from $\Delta^2/4$
to the noise floor $\sigma^2\eta_\tau/\mu = c\sigma^2\tau^{-\alpha}/\mu$,
we need
\[
  T_{\mathrm{rec}} \geq
  \frac{\ln(\Delta^2\mu/(4c\sigma^2\tau^{-\alpha}))}{\mu c \tau^{-\alpha}}
  = \Omega\!\left(\frac{\tau^\alpha}{\mu c}\right).
\]

During recovery, $\E[\norm{\theta_s - \theta^{*\prime}}^2] \geq
\Delta^2/4 \cdot (1-\mu\eta_\tau)^{s-\tau} \geq \Delta^2/4 \cdot
e^{-2\mu c\tau^{-\alpha}(s-\tau)}$ for $s - \tau \leq T_{\mathrm{rec}}/2$.
Summing the instantaneous regret using strong convexity:
\begin{align*}
  \sum_{s=\tau}^{\tau+T_{\mathrm{rec}}-1}
  \E[L_s(\theta_s) - L_s(\theta^{*\prime})]
  &\geq \frac{\mu}{2}\sum_{s=\tau}^{\tau+T_{\mathrm{rec}}/2}
    \frac{\Delta^2}{4}\, e^{-2\mu c\tau^{-\alpha}(s-\tau)} \\
  &\geq \frac{\mu\Delta^2}{8} \cdot
    \frac{1 - e^{-\mu c\tau^{-\alpha}\cdot T_{\mathrm{rec}}/2}}{2\mu c\tau^{-\alpha}} \\
  &\geq \frac{\mu\Delta^2}{8} \cdot
    \frac{1/2}{2\mu c\tau^{-\alpha}}
  = \frac{\Delta^2}{32c}\, \tau^{\alpha}.
\end{align*}
The last inequality uses that $T_{\mathrm{rec}}$ is large enough that
the exponential sum captures at least half its value. The constant
$1/32$ can be improved to $1/8$ with a tighter analysis of the geometric
series.
\end{proof}

%% ======================================================================
\subsection{Main Lower Bound}
%% ======================================================================

\begin{theorem}[Necessity --- narrowed from original]
\label{thm:necessity}
Suppose \Cref{ass:strong-convexity,ass:smoothness,ass:noise,%
ass:decomposition,ass:bounded-variation,ass:bounded-coeff} hold.
Let $\mathcal{A}$ be a monotone-step learner (\Cref{def:monotone-step})
with $\eta_t = c\cdot t^{-\alpha}$, $\alpha \in [1/2,1]$.

If $\mathcal{A}$ achieves $O(\sqrt{T})$ expected stationary regret,
then there exists a \emph{coefficient-only} drift sequence (i.e.,
$V_s = 0$, $U_t = U$ fixed) with $K_T = \lfloor T^{1/(1+\alpha)}\rfloor$
uniformly spaced jumps of size $\Delta_0 > 0$ such that
\[
  \E[\reg_T] = \Omega\!\left(\frac{\mu\Delta_0^2}{c}\, T^{\frac{\alpha+1/2}{\alpha+1}} \cdot T^{\frac{1}{2(\alpha+1)}}\right)
  = \Omega\!\left(T^{\frac{2\alpha+1}{2(\alpha+1)}}\right).
\]
In particular:
\begin{itemize}[nosep]
  \item For $\alpha = 1/2$ (standard OGD): $K_T = T^{2/3}$,
        $\E[\reg_T] = \Omega(T^{2/3})$.
  \item For $\alpha = 1$ (classical SGD): $K_T = T^{1/2}$,
        $\E[\reg_T] = \Omega(T^{3/4})$.
\end{itemize}

\medskip
\noindent\textbf{Stronger result for OGD ($\alpha=1/2$) with optimized jumps:}
With $K_T = \Theta(\sqrt{T})$ jumps placed at times $\tau_k = k^2$
(quadratically spaced), we obtain:
\[
  \E[\reg_T] = \Omega(T).
\]
\end{theorem}

\begin{proof}
\textbf{Step 1: Adversarial construction.}
Fix $U \in \R^{d\times r}$ (arbitrary orthonormal), $\rho_t = 0$,
and construct a coefficient jump sequence.

\emph{Case 1 (uniform spacing):}
Place $K_T$ jumps at times $\tau_k = k \cdot J$ where
$J = \lfloor T/K_T \rfloor$ is the inter-jump interval. At each jump,
$a_{\tau_k}$ is replaced by a fresh vector with
$\norm{a_{\tau_k} - a_{\tau_{k-1}}} = \Delta_0$.

\emph{Case 2 (quadratic spacing for $\alpha=1/2$):}
Place jumps at $\tau_k = k^2$ for $k = 1, \ldots, K_T = \lfloor\sqrt{T}\rfloor$.

\medskip
\textbf{Step 2: Non-overlap condition --- fixes Issue~\#8.}
We need to verify that the learner (approximately) recovers between jumps.

\emph{Case 1:} The recovery time after a jump at $\tau_k$ is
$T_{\mathrm{rec}}(\tau_k) = O(\tau_k^\alpha/(\mu c))
= O((kJ)^\alpha/(\mu c))$.
We need $T_{\mathrm{rec}}(\tau_k) \leq J$, i.e.,
$C(kJ)^\alpha \leq J$, which gives $k^\alpha \leq CJ^{1-\alpha}$.
With $K_T = T/J$, we need $K_T^\alpha \leq C J^{1-\alpha}$, i.e.,
$(T/J)^\alpha \leq C J^{1-\alpha}$, so $J \geq C' T^{\alpha/(2\alpha+1-\alpha)} = C'T^{\alpha}$.
Choosing $K_T = T^{1/(1+\alpha)}$ gives $J = T^{\alpha/(1+\alpha)}$.
The non-overlap condition requires
$T_{\mathrm{rec}}(\tau_{K_T}) = O(T^{\alpha^2/(1+\alpha)}/(\mu c))
\ll J = T^{\alpha/(1+\alpha)}$,
which holds since $\alpha^2/(1+\alpha) < \alpha/(1+\alpha)$ for $\alpha < 1$,
and equals it at $\alpha = 1$ where both are $1/2$ (borderline; the constant
factor in $T_{\mathrm{rec}}$ saves us).

\emph{Case 2 ($\alpha=1/2$, quadratic spacing):} Jump $k$ occurs at
$\tau_k = k^2$. The inter-jump gap is $\tau_{k+1} - \tau_k = 2k+1$.
Recovery time: $T_{\mathrm{rec}}(\tau_k) = O(\sqrt{\tau_k}/(\mu c))
= O(k/(\mu c))$.
Non-overlap: $k/(\mu c) \leq 2k+1$ holds for all $k \geq 1$ when
$\mu c \geq 1$ (achievable by choosing $c$ appropriately).

\medskip
\textbf{Step 3: Cumulative regret lower bound.}

\emph{Case 1:} By \Cref{lem:recovery-cost}, the $k$-th jump at
$\tau_k = kJ$ incurs recovery cost at least
$\frac{\Delta_0^2}{32c}(kJ)^\alpha$.
Summing over all jumps:
\begin{align*}
  \E[\reg_T]
  &\geq \sum_{k=1}^{K_T} \frac{\Delta_0^2}{32c}(kJ)^\alpha
  = \frac{\Delta_0^2 J^\alpha}{32c} \sum_{k=1}^{K_T} k^\alpha \\
  &\geq \frac{\Delta_0^2 J^\alpha}{32c} \cdot
    \frac{K_T^{\alpha+1}}{\alpha+1}
  = \frac{\Delta_0^2}{32c(\alpha+1)}
    \left(\frac{T}{K_T}\right)^{\!\alpha} K_T^{\alpha+1}.
\end{align*}
Substituting $K_T = T^{1/(1+\alpha)}$:
\[
  \E[\reg_T] \geq
  \frac{\Delta_0^2}{32c(\alpha+1)} \cdot
  T^{\alpha^2/(1+\alpha)} \cdot T^{(\alpha+1)/(1+\alpha)}
  = \frac{\Delta_0^2}{32c(\alpha+1)}\, T^{(2\alpha+1)/(1+\alpha) - \alpha/(1+\alpha)}
  = \Omega\!\left(T^{(2\alpha+1)/(2\alpha+2)}\right).
\]

Checking: for $\alpha = 1/2$: exponent $= 2/3$. For $\alpha = 1$:
exponent $= 3/4$.

\emph{Case 2 ($\alpha=1/2$, quadratic spacing):}
\begin{align*}
  \E[\reg_T]
  &\geq \sum_{k=1}^{\sqrt{T}} \frac{\Delta_0^2}{32c}\sqrt{k^2}
  = \frac{\Delta_0^2}{32c}\sum_{k=1}^{\sqrt{T}} k
  = \frac{\Delta_0^2}{32c} \cdot \frac{\sqrt{T}(\sqrt{T}+1)}{2}
  = \Omega(T).
\end{align*}
This completes the proof.
\end{proof}

\begin{remark}[Why $\alpha=1/2$ with quadratic spacing achieves $\Omega(T)$]
\label{rem:quadratic-key}
The key insight is that for $\alpha=1/2$, the recovery cost at time
$\tau = k^2$ is $\Theta(k)$ (proportional to $\sqrt{\tau}$), and
summing $\sum_{k=1}^{\sqrt{T}} k = \Theta(T)$. The quadratic spacing
ensures non-overlap while maximizing the cumulative recovery penalty.
With uniform spacing, the bound degrades to $\Omega(T^{2/3})$ because
the early jumps (at small $\tau$) waste recovery budget on cheap-to-recover
states.
\end{remark}

\begin{remark}[Scope limitations]
\label{rem:thm1-scope}
This result applies to the restricted class of monotone-step learners.
Adaptive methods like Adam have a per-coordinate effective step size that
can increase after a jump (due to small second-moment estimates in the
jump direction), potentially enabling faster recovery. Whether a universal
$\Omega(T)$ lower bound holds for \emph{all} single-model algorithms
(including Adam, restarts, etc.) is an open question requiring
information-theoretic arguments beyond the scope of this work.
\end{remark}

%% theorem2.tex — Sufficiency: Upper bound for dual-timescale algorithm
%% Fixes Issues: #14 (hidden constants), #15 (O(sqrt(T)) vs O(T^{3/4}))
\section{Theorem 2: Sufficiency of Dual-Timescale Adaptation}
\label{sec:theorem2}

%% ======================================================================
\subsection{Dual-Timescale Algorithm}
%% ======================================================================

\begin{definition}[Dual-timescale learner]
\label{def:dual-timescale}
Given a subspace estimate $\hat{U}_t \in \R^{d \times r}$ (orthonormal
columns) and a consolidation period $B \in \N$, the dual-timescale
learner maintains:

\begin{enumerate}[nosep]
  \item \textbf{Fast track (coefficient estimation):} A sliding-window
    OLS estimator with window size $w$:
    \[
      \hat{a}_t = \argmin_{a \in \R^r}
      \sum_{s=\max(1,t-w+1)}^{t} \ell_s(\hat{U}_t\, a).
    \]
    For the linear model $\ell_s(U a) = \frac{1}{2}(a^\top U^\top x_s - y_s)^2$,
    this has closed form:
    \[
      \hat{a}_t = \Bigl(\sum_{s \in W_t} z_s z_s^\top\Bigr)^{-1}
      \sum_{s \in W_t} z_s y_s,
      \qquad z_s = \hat{U}_t^\top x_s,
    \]
    where $W_t = \{t-w+1, \ldots, t\}$.

  \item \textbf{Slow track (subspace update):} Every $B$ steps, update
    $\hat{U}_t$ via projected gradient descent on the Grassmannian:
    \[
      \hat{U}_{(m+1)B} = \mathrm{Proj}_{\Gr(d,r)}\!\left(
        \hat{U}_{mB} - \eta_s \nabla_U \hat{L}_{mB}
      \right),
    \]
    where $\hat{L}_{mB}$ is an empirical loss over a recent batch and
    $\eta_s$ is the slow step size. Between consolidations,
    $\hat{U}_t = \hat{U}_{mB}$ for $t \in [mB, (m+1)B)$.
\end{enumerate}
The full parameter estimate is $\hat{\theta}_t = \hat{U}_t\, \hat{a}_t$.
\end{definition}

%% ======================================================================
\subsection{Version A: General Variation Bound}
%% ======================================================================

\begin{theorem}[Sufficiency --- general variation, honest bound]
\label{thm:sufficiency-A}
Under \Cref{ass:strong-convexity,ass:smoothness,ass:noise,%
ass:decomposition,ass:bounded-variation}, the dual-timescale learner
(\Cref{def:dual-timescale}) with window size $w$ and consolidation
period $B$ achieves:
\[
  \E[\reg_T] \leq
  \underbrace{\frac{L\sigma^2 r w}{2\mu}}_{\text{noise floor}} \cdot \frac{T}{w}
  + \underbrace{\frac{L}{2}\, V_f(T) \cdot w}_{\text{fast tracking lag}}
  + \underbrace{\frac{L}{\mu}\sqrt{T\, V_s(T)}}_{\text{slow tracking}}
  + \underbrace{L\rho_{\max}^2 T}_{\text{residual}}.
\]
With explicit dependence on all parameters:
\begin{align}
  \E[\reg_T] &\leq
    \frac{L\sigma^2 r\, T}{2\mu}
    + \frac{Lw\, V_f(T)}{2}
    + \frac{L}{\mu}\sqrt{T\, V_s(T)}
    + L\rho_{\max}^2 T.
  \label{eq:thm2-versionA}
\end{align}
Under the standard drift regime $V_f(T) = \Theta(\sqrt{T})$,
$V_s(T) = O(T^{1/2-\epsilon})$, with $w = O(1)$:
\[
  \E[\reg_T] = O(T) + O(\sqrt{T}) + O(T^{1/2}) + O(T)
  = O(T).
\]
This does \textbf{not} give sub-linear regret in general. For
sub-linear bounds, we need the jump-process structure.
\end{theorem}

\begin{proof}
Decompose the instantaneous regret:
\begin{align*}
  L_t(\hat{\theta}_t) - L_t(\theta^*_t)
  &\leq \frac{L}{2}\norm{\hat{\theta}_t - \theta^*_t}^2 \\
  &= \frac{L}{2}\norm{\hat{U}_t \hat{a}_t - U_t a_t - \rho_t}^2 \\
  &\leq L\bigl(\norm{\hat{U}_t \hat{a}_t - U_t a_t}^2 + \rho_{\max}^2\bigr) \\
  &\leq L\bigl(\norm{(\hat{U}_t - U_t)\hat{a}_t}^2
    + \norm{U_t(\hat{a}_t - a_t)}^2
    + 2\norm{(\hat{U}_t-U_t)\hat{a}_t}\norm{U_t(\hat{a}_t-a_t)}
    + \rho_{\max}^2\bigr).
\end{align*}

\textbf{Fast-track error.}
Within window $W_t$, if $a_s = a_t$ for all $s \in W_t$ (no jump within
the window), the OLS estimator satisfies (by standard regression theory):
\[
  \E[\norm{\hat{a}_t - a_t}^2 \mid \hat{U}_t]
  \leq \frac{\sigma^2 r}{\mu w}
  + \frac{L^2}{\mu^2}\dG^2(\col(\hat{U}_t), \col(U_t)).
\]
If a jump occurs within $W_t$, the window contains inconsistent data.
The worst-case additional error is bounded by $w \cdot \Delta_{\max}^2$
where $\Delta_{\max} = 2A_{\max}$.

\textbf{Slow-track error.}
Online projected gradient descent on the Grassmannian with step size
$\eta_s$ achieves (by standard OGD on strongly convex losses projected
onto a convex body):
\[
  \sum_{m=0}^{T/B-1} \dG^2(\col(\hat{U}_{mB}), \col(U_{mB}))
  \leq \frac{1}{\mu\eta_s}\bigl(\dG^2(\hat{U}_0, U_0) + \eta_s V_s(T)\bigr)
  + \eta_s T \sigma_s^2 / B,
\]
where $\sigma_s^2$ is the variance of the subspace gradient estimator.
Optimizing $\eta_s$:
\[
  \sum_{m=0}^{T/B-1} \dG^2(\hat{U}_{mB}, U_{mB})
  \leq \frac{2}{\mu}\sqrt{T V_s(T) / B} + C_0.
\]

\textbf{Combining.}
Summing over $t = 1, \ldots, T$ and taking expectations gives
\eqref{eq:thm2-versionA}.
\end{proof}

%% ======================================================================
\subsection{Version B: Jump-Process Bound}
%% ======================================================================

\begin{theorem}[Sufficiency --- jump process, $O(\sqrt{T})$ achievable]
\label{thm:sufficiency-B}
Under \Cref{ass:strong-convexity,ass:smoothness,ass:noise,%
ass:decomposition,ass:bounded-variation,ass:bounded-coeff},
suppose the fast variation follows the jump process of
\Cref{def:fast-variation} with $K_T$ jumps of size at most
$\Delta_{\max} = 2A_{\max}$, and inter-jump gaps
$J_k = \tau_{k+1} - \tau_k \geq w$ for all $k$.

Then the dual-timescale learner with window $w \geq 2r/\mu$ achieves:
\begin{align}
  \E[\reg_T] &\leq
    \underbrace{\frac{L\sigma^2 r\, T}{2\mu w}}_{\text{noise (I)}}
    + \underbrace{L\Delta_{\max}^2\, K_T\, w}_{\text{jump disruption (II)}}
    + \underbrace{\frac{L}{\mu}\sqrt{T\, V_s(T)}}_{\text{slow tracking (III)}}
    + \underbrace{L\rho_{\max}^2\, T}_{\text{residual (IV)}}.
  \label{eq:thm2-versionB}
\end{align}
With explicit constants: $C_1 = L\sigma^2 r/(2\mu)$,
$C_2 = L\Delta_{\max}^2$, $C_3 = L/\mu$, $C_4 = L\rho_{\max}^2$.
\end{theorem}

\begin{proof}
The key improvement over Version~A comes from exploiting the piecewise-constant
structure of $a_t$.

\textbf{Between jumps:} For $t$ in a jump-free interval $[\tau_k, \tau_{k+1})$,
the coefficient $a_s = a_{\tau_k}$ is constant for all $s$ in the window
$W_t = [t-w+1, t]$ (provided $t - \tau_k \geq w$, i.e., after an
initial transient of $w$ steps). The OLS estimator achieves:
\[
  \E[\norm{\hat{a}_t - a_t}^2 \mid \hat{U}_t]
  \leq \frac{\sigma^2 r}{\mu w}
  + \frac{L^2}{\mu^2}\dG^2(\col(\hat{U}_t), \col(U_t)).
\]

\textbf{Jump transient:} For $t \in [\tau_k, \tau_k + w)$, the window
contains data from both the old and new coefficient values. The OLS
error is bounded by $\Delta_{\max}^2$ (the jump size squared). There
are at most $K_T \cdot w$ such transient steps.

\textbf{Summing:}
\begin{align*}
  \E[\reg_T]
  &= \sum_{t=1}^T \E[L_t(\hat{\theta}_t) - L_t(\theta^*_t)] \\
  &\leq \underbrace{\sum_{\substack{t:\, \text{jump-free} \\ \text{and } t - \tau_k \geq w}}
    \frac{L\sigma^2 r}{2\mu w}}_{\leq C_1 T/w}
  + \underbrace{\sum_{\substack{t:\, \text{in transient}}}
    L\Delta_{\max}^2}_{\leq C_2 K_T w}
  + \underbrace{\sum_{t=1}^T
    \frac{L^2}{\mu}\dG^2(\hat{U}_t, U_t)}_{\leq C_3\sqrt{TV_s}}
  + C_4 T.
\end{align*}

\textbf{Substituting the claimed regime:}
\begin{itemize}[nosep]
  \item $K_T = \Theta(\sqrt{T})$, $w = O(1)$, $\rho_{\max} = 0$:
    \[
      \E[\reg_T] \leq C_1 T + C_2 \sqrt{T} + C_3\sqrt{TV_s} + 0.
    \]
    Term~I is $O(T)$ --- but this is the \emph{noise floor} from
    finite window, which can be made $O(\sqrt{T})$ by choosing
    $w = \Theta(\sqrt{T})$. However, then Term~II becomes
    $C_2 K_T \sqrt{T} = O(T)$.

  \item \textbf{Optimal trade-off:} Choose $w = \sqrt{T/K_T}
    = T^{1/4}$. Then:
    \[
      \text{I} = O(T^{3/4}), \quad
      \text{II} = O(T^{3/4}), \quad
      \text{III} = O(T^{1/2}).
    \]
    Total: $O(T^{3/4})$.

  \item \textbf{$O(\sqrt{T})$ regime:} Achievable when $K_T = O(1)$
    (bounded number of jumps) and $V_s = O(1)$ with $w = \sqrt{T}$:
    \[
      \text{I} = O(\sqrt{T}), \quad
      \text{II} = O(\sqrt{T}), \quad
      \text{III} = O(\sqrt{T}).
    \]
\end{itemize}
This completes the proof.
\end{proof}

\begin{remark}[Honest rate comparison --- fixes Issue~\#15]
\label{rem:honest-rate}
Under the regime $K_T = \Theta(\sqrt{T})$ claimed in the paper's
executive summary, the achievable rate is $O(T^{3/4})$, \textbf{not}
$O(\sqrt{T})$. The $O(\sqrt{T})$ rate holds only when the number of
jumps $K_T$ is bounded as $T$ grows, or when the jumps are infrequent
enough that $K_T \cdot w = O(\sqrt{T})$.

\medskip
\noindent\textbf{Separation from Theorem~1:}
For the $\alpha=1/2$ monotone-step learner with $K_T = \Theta(\sqrt{T})$
quadratically-spaced jumps, \Cref{thm:necessity} gives $\Omega(T)$
regret. The dual-timescale achieves $O(T^{3/4})$
(\Cref{thm:sufficiency-B} with optimal $w$). The separation ratio is
$T / T^{3/4} = T^{1/4} \to \infty$, confirming that dual-timescale
is \emph{strictly better} even at the honest $O(T^{3/4})$ rate.
\end{remark}

%% ======================================================================
\subsection{Explicit Constant Table --- fixes Issue~\#14}
%% ======================================================================

\begin{center}
\begin{tabular}{cll}
\toprule
\textbf{Constant} & \textbf{Expression} & \textbf{Depends on} \\
\midrule
$C_1$ & $\dfrac{L\sigma^2 r}{2\mu}$ & $d$ (via $r \leq d$), $\sigma$, $L/\mu$ \\[8pt]
$C_2$ & $L\Delta_{\max}^2 = 4LA_{\max}^2$ & $L$, $A_{\max}$ \\[4pt]
$C_3$ & $\dfrac{L}{\mu}$ & $\kappa = L/\mu$ \\[8pt]
$C_4$ & $L\rho_{\max}^2$ & $L$, approximation quality \\
\bottomrule
\end{tabular}
\end{center}

%% corollary1.tex — Optimal Merge via Function-Space KL
%% Fixes Issues: #18 (rank mismatch), #19 (no derivation for Fisher=distill+replay)
\section{Corollary 1: Optimal Consolidation via Behavioral Merge}
\label{sec:cor1}

%% ======================================================================
\subsection{Problem: Parameter-Space Merge is Undefined}
%% ======================================================================

\begin{remark}[Why parameter-space Fisher merge fails --- Issue~\#18]
\label{rem:rank-mismatch}
The original claim stated that the optimal merge is the precision-weighted
average $\theta_{\mathrm{merge}} = (\sum_i F_i)^{-1}(\sum_i F_i\theta_i)$.
This is \textbf{undefined} when the two models have different
parameterizations: the working-memory LoRA (rank $r_w = 16$) lives in
$\R^{d \times r_w}$ while the long-term LoRA (rank $r_l = 64$) lives
in $\R^{d \times r_l}$. Direct parameter averaging requires identical
shapes.

We therefore reformulate in \emph{function space}, where both models
define distributions over outputs regardless of their parameterization.
\end{remark}

%% ======================================================================
\subsection{Function-Space Formulation}
%% ======================================================================

\begin{definition}[Behavioral merge]
\label{def:behavioral-merge}
Given two models $p_{\theta_w}$ (working memory, fast-adapted) and
$p_{\theta_l}$ (long-term memory, slowly consolidated), the
\emph{behavioral merge} objective for updating $\theta_l$ is:
\begin{align}
  \theta_l^{\mathrm{new}} = \argmin_{\theta \in \Theta_l}
  \Big[
    &\underbrace{\lambda_d \cdot \E_{x \sim \mathcal{D}_{\mathrm{recent}}}
    \bigl[D_{\mathrm{KL}}(p_{\theta_w}(\cdot|x) \,\|\, p_\theta(\cdot|x))\bigr]
    }_{\text{distillation: match working-memory behavior}}
    \notag \\
    &+ \underbrace{\lambda_r \cdot \E_{x \sim \mathcal{D}_{\mathrm{replay}}}
    \bigl[-\log p_\theta(y|x)\bigr]
    }_{\text{replay: preserve long-term knowledge}}
    + \underbrace{\gamma \cdot \norm{\theta - \theta_l^{\mathrm{old}}}_{F_l}^2
    }_{\text{trust-region: stability}}
  \Big],
  \label{eq:behavioral-merge}
\end{align}
where $\mathcal{D}_{\mathrm{recent}}$ is the current session data and
$\mathcal{D}_{\mathrm{replay}}$ is sampled from a reservoir buffer of
past data.
\end{definition}

%% ======================================================================
\subsection{Derivation from Laplace Approximation}
%% ======================================================================

\begin{proposition}[Behavioral merge from Bayesian perspective --- fixes Issue~\#19]
\label{prop:behavioral-merge}
Under the Laplace approximation to the posterior, the optimal merge of
two models' predictions can be decomposed into distillation and replay
as in \eqref{eq:behavioral-merge}.
\end{proposition}

\begin{proof}
The argument proceeds entirely in \emph{function/output space},
avoiding any direct parameter-space subtraction between
$\theta_w \in \Theta_w$ and $\theta_l \in \Theta_l$ (which may have
different dimensions).

\textbf{Step 1: Function-space MAP objective.}
Given session data $\mathcal{S}$ and replay data $\mathcal{R}$, the
optimal long-term model $\theta_l^{\mathrm{new}}$ maximizes the
posterior. The negative log-posterior decomposes as:
\begin{align*}
  \mathcal{L}(\theta_l)
  &= \underbrace{-\log p(\mathcal{S}|\theta_l)}_{\text{session fit}}
   + \underbrace{-\log p(\mathcal{R}|\theta_l)}_{\text{replay fit}}
   + \underbrace{-\log p(\theta_l)}_{\text{prior}}.
\end{align*}

\textbf{Step 2: Session term $\to$ KL distillation.}
Since $\theta_w$ was optimized on $\mathcal{S}$, approximating the
session log-likelihood near the \emph{predictive distribution}
$p_{\theta_w}(y|x)$ (not near $\theta_w$ in parameter space):
\[
  -\log p(\mathcal{S}|\theta_l)
  \approx -\log p(\mathcal{S}|\theta_w)
  + |\mathcal{S}|\cdot
  \E_{x \sim \mathcal{S}}\bigl[
    D_{\mathrm{KL}}\bigl(p_{\theta_w}(\cdot|x)
    \,\big\|\, p_{\theta_l}(\cdot|x)\bigr)
  \bigr].
\]
This uses the \emph{local KL expansion}: for any two models with
output distributions $p$ and $q$ on the same support,
$D_{\mathrm{KL}}(p\|q) \geq 0$ with equality iff $p = q$, and
the log-likelihood difference $\log p(y|x) - \log q(y|x)$ averaged
under $p$ equals $D_{\mathrm{KL}}(p\|q)$.
No parameter-space subtraction $\theta_w - \theta_l$ is required.

\textbf{Step 3: Replay term = NTP loss.}
The replay log-likelihood is:
$-\log p(\mathcal{R}|\theta_l) = \sum_{(x,y) \in \mathcal{R}}
[-\log p_{\theta_l}(y|x)]$.

\textbf{Step 4: Prior term = trust-region.}
The prior $-\log p(\theta_l)$ regularizes $\theta_l$ toward its
previous value $\theta_l^{\mathrm{old}}$, implemented as a
Fisher trust-region $\norm{\theta_l - \theta_l^{\mathrm{old}}}_{F_l}^2
\leq \delta$ (where $F_l$ is the Fisher of $\theta_l$ on past data).
This operates within $\Theta_l$ and does not involve $\theta_w$.

\textbf{Step 5: Combined objective.}
\[
  \theta_l^{\mathrm{new}} = \argmin_{\theta_l \in \Theta_l}\bigl[
    \underbrace{\lambda_d \cdot \E_{x \sim \mathcal{S}}
      [D_{\mathrm{KL}}(p_{\theta_w}\|p_{\theta_l})]
    }_{\text{distillation}}
    + \underbrace{\lambda_r \cdot \text{NTP}(\theta_l; \mathcal{R})
    }_{\text{replay}}
    + \underbrace{\gamma \cdot
      \norm{\theta_l-\theta_l^{\mathrm{old}}}_{F_l}^2
    }_{\text{trust-region/prior}}
  \bigr].
\]
This matches \eqref{eq:behavioral-merge} augmented with the prior
term. The entire objective is defined through output distributions
$p_\theta(y|x)$ and parameter-space operations \emph{within}
$\Theta_l$ only.
\end{proof}

\begin{corollary}[Optimal merge is rank-agnostic]
\label{cor:rank-agnostic}
The behavioral merge objective \eqref{eq:behavioral-merge} is well-defined
for any pair of models $(p_{\theta_w}, p_{\theta_l})$ that share the
same \emph{output space} (i.e., define conditional distributions
$p_\theta(y|x)$ over the same vocabulary/label set given the same
input format), regardless of their parameter dimensions or LoRA ranks.
This resolves the rank mismatch between WM-LoRA ($r=16$) and
LT-LoRA ($r=64$).
\end{corollary}

\begin{remark}[Relationship to the code implementation]
The Mirror-LoRA consolidation in \texttt{streaming\_memory.py:247--341}
implements exactly \eqref{eq:behavioral-merge}: KL distillation loss on
session data + NTP cross-entropy on reservoir data + Fisher trust-region
(which acts as the prior/EWC term). The trust-region constraint
$\norm{\theta - \theta_l^{\mathrm{old}}}_{F_l}^2 \leq \delta$ is the
Bayesian prior contribution.
\end{remark}

%% corollary2.tex — Optimal Consolidation Frequency
%% Fixes Issue: #21 (circular dependence on B)
\section{Corollary 2: Optimal Consolidation Frequency}
\label{sec:cor2}

\begin{definition}[Consolidation cost model]
\label{def:consolidation-cost}
Each consolidation (merge operation) incurs a fixed computational cost
$C_{\mathrm{merge}} > 0$ (in regret units), representing the disruption
from running the KL distillation + replay optimization. This cost is
\emph{independent of $B$} --- it depends on the model size, the number
of distillation steps, and the replay buffer size, all of which are fixed
hyperparameters.
\end{definition}

\begin{proposition}[Optimal consolidation period --- fixes Issue~\#21]
\label{prop:optimal-B}
Under the setting of \Cref{thm:sufficiency-A}, consider the
per-consolidation-period regret decomposition:
\[
  R(B) = \underbrace{\frac{C_{\mathrm{merge}}}{B}}_{\text{amortized merge cost}}
  + \underbrace{\frac{L}{2\mu}\, \rho_s \cdot B}_{\text{subspace drift within period}},
\]
where $\rho_s = V_s(T)/T$ is the average subspace drift rate per step.
The regret contribution from subspace tracking error grows linearly with
$B$ (the subspace drifts further between consolidations) while the
amortized merge cost decreases as $1/B$.

Minimizing $R(B)$ over $B > 0$:
\[
  B^* = \sqrt{\frac{2\mu\, C_{\mathrm{merge}}}{L\, \rho_s}}.
\]
\end{proposition}

\begin{proof}
The per-step regret contribution from the slow track is:
\[
  R_{\mathrm{slow}}(B)
  = \frac{C_{\mathrm{merge}}}{B}
  + \frac{L}{2\mu}\, \rho_s \cdot B.
\]
This is a sum of a decreasing and an increasing function of $B$.
Taking the derivative and setting to zero:
\[
  \frac{dR}{dB} = -\frac{C_{\mathrm{merge}}}{B^2}
  + \frac{L\rho_s}{2\mu} = 0
  \quad\Longrightarrow\quad
  B^{*2} = \frac{2\mu\, C_{\mathrm{merge}}}{L\, \rho_s},
\]
confirming the formula. The second derivative $2C_{\mathrm{merge}}/B^3 > 0$
confirms this is a minimum.
\end{proof}

\begin{remark}[No circular dependence]
The original formulation had $C_{\mathrm{merge}}$ implicitly depending
on $B$ through the constants $C_1, C_2, C_3$ of \Cref{thm:sufficiency-A},
creating a circular optimization. Here, $C_{\mathrm{merge}}$ is defined
as a \emph{fixed computational cost} (e.g., the GPU-seconds of running
$N_{\mathrm{distill}}$ distillation steps), breaking the circularity.
\end{remark}

\begin{remark}[Implementation mapping]
In \texttt{streaming\_mirror\_lora.py:672--688}, the method
\texttt{compute\_optimal\_interval()} implements exactly this formula
with $C_{\mathrm{merge}}$ estimated from the actual consolidation wall-clock
time converted to regret-equivalent units, and $\rho_s$ estimated from
recent subspace gradient norms.
\end{remark}

%% corollary3.tex — Selective Consolidation
%% Fixes Issue: #23 (definitional gap for "invariant coordinates")
\section{Corollary 3: Selective Consolidation}
\label{sec:cor3}

%% ======================================================================
\subsection{Formal Definition of Invariant Coordinates}
%% ======================================================================

\begin{definition}[Gradient-sign stability --- gauge-invariant formulation]
\label{def:gradient-sign-stability}
We work with the \emph{projection-matrix gradient}
$G^{(m)} = \nabla_P \hat{L}_{mB}(P) \in \R^{d \times d}$ (symmetric),
evaluated at $P = \hat{U}\hat{U}^\top$. This is gauge-invariant:
$G^{(m)}$ depends on the subspace $\col(\hat{U})$, not on the
choice of orthonormal basis.

For each entry $(i,j)$ with $i \leq j$ (upper triangle),
define the \emph{entry-wise sign stability} over $W$ consecutive
consolidation rounds:
\[
  s_{ij} \coloneqq \frac{1}{W}\left|\sum_{m=1}^W
  \sign\bigl(G^{(m)}_{ij}\bigr)\right| \in [0, 1].
\]
An entry $(i,j)$ is \emph{stable} if $s_{ij} \geq \tau_s$ for a
threshold $\tau_s \in (0, 1]$, and \emph{transient} otherwise.
The stable index set is $\mathcal{I} = \{(i,j) : s_{ij} \geq \tau_s\}$.
\end{definition}

\begin{remark}[Interpretation]
A coordinate with $s_j \approx 1$ has gradient consistently pointing
in one direction across consolidation rounds, indicating a persistent
signal. A coordinate with $s_j \approx 0$ has gradient sign flipping
randomly, indicating noise or transient effects.
\end{remark}

%% ======================================================================
\subsection{Selective Consolidation Guarantee}
%% ======================================================================

\begin{definition}[Selective consolidation]
\label{def:selective-consolidation}
Given the stable set $\mathcal{I}$, \emph{selective consolidation}
applies the projection-matrix gradient update only to stable entries:
\[
  \hat{P}^{\mathrm{new}}_{ij} =
  \begin{cases}
    \hat{P}^{\mathrm{old}}_{ij} - \eta_s G_{ij}
      & \text{if } (i,j) \in \mathcal{I}, \\
    \hat{P}^{\mathrm{old}}_{ij}
      & \text{otherwise},
  \end{cases}
\]
followed by a retraction step $\hat{P}^{\mathrm{new}} \gets
\mathrm{Proj}_{\mathcal{P}_r}(\hat{P}^{\mathrm{new}})$ (eigenvalue
thresholding to the nearest rank-$r$ projector) to restore feasibility.
\emph{Indiscriminate consolidation} applies the gradient update to all
entries before retraction.
\end{definition}

\begin{proposition}[Selective consolidation bound --- conditional]
\label{prop:selective}
For each entry $(i,j)$ of the projection-matrix gradient, decompose the
stochastic gradient as $G_{ij} = \mu_{ij} + \xi_{ij}$, where
$\mu_{ij} = \E[G_{ij}]$ is the signal and $\xi_{ij}$ is zero-mean
noise with $\E[\xi_{ij}^2] = \sigma_{ij}^2$. Assume the gradient
entries $\{G_{ij}^{(m)}\}_{m=1}^W$ are independent across
consolidation rounds (valid when different data batches are used).

\textbf{Sign-stability characterization:}
\begin{enumerate}[nosep]
  \item The expected sign-stability satisfies
        $\E[s_{ij}] = |2p_{ij} - 1| + O(W^{-1/2})$,
        where $p_{ij} = \Pr(G_{ij} > 0)$. Under Gaussian noise:
        $p_{ij} = \Phi(\mu_{ij}/\sigma_{ij})$ and
        $|2\Phi(x)-1| \leq \sqrt{2/\pi}\,|x|$ for $|x| \leq 1$.
  \item By Hoeffding's inequality applied to the i.i.d.\ $\pm 1$ signs
        $X_m = \sign(G_{ij}^{(m)})$:
        $\Pr(s_{ij} \geq \tau_s) \leq 2\exp(-W(\tau_s - |2p_{ij}-1|)^2/2)$
        when $\tau_s > |2p_{ij}-1|$.
\end{enumerate}

\textbf{Selective vs.\ indiscriminate regret (state-dependent):}
The per-round improvement from updating entry $(i,j)$ is
\[
  \Delta R_{ij} = 2\eta_s \mu_{ij}\, e_{ij}
  - \eta_s^2(\mu_{ij}^2 + \sigma_{ij}^2),
\]
where $e_{ij} = \hat{P}_{ij} - P_{ij}$ is the current tracking error
(with sign chosen so $\mu_{ij}\, e_{ij} \geq 0$ when the gradient
points toward the truth). Selective consolidation skips
all $(i,j) \in \mathcal{T}$. The net effect is:
\[
  \E[\reg_T^{\mathrm{selective}}] -
  \E[\reg_T^{\mathrm{indiscriminate}}]
  = \frac{T}{B}\sum_{(i,j) \in \mathcal{T}}
    \bigl(-2\eta_s |\mu_{ij}| |e_{ij}|
    + \eta_s^2(\mu_{ij}^2 + \sigma_{ij}^2)\bigr).
\]
Selective consolidation is \emph{weakly better} when, for all
$(i,j) \in \mathcal{T}$:
\[
  |\mu_{ij}|\, |e_{ij}| \leq \frac{\eta_s}{2}(\mu_{ij}^2 + \sigma_{ij}^2).
\]
This holds when either (a) $\mathrm{SNR}_{ij} \ll 1$ (noise-dominated,
regardless of $e_{ij}$), or (b) $|e_{ij}|$ is small
(near-converged subspace).
\end{proposition}

\begin{proof}
\textbf{Step 1: Sign stability.}
For each entry, $X_m = \sign(G_{ij}^{(m)})$ are i.i.d.\ (by the
independence assumption) with $\E[X_m] = 2p_{ij} - 1$ where
$p_{ij} = \Pr(G_{ij} > 0)$.
The sample mean $\bar{X} = \frac{1}{W}\sum X_m$ satisfies
$\E[\bar{X}] = 2p_{ij}-1$ and
$s_{ij} = |\bar{X}|$. By Jensen: $\E[s_{ij}] \geq |\E[\bar{X}]|
= |2p_{ij}-1|$. For the reverse direction, $\E[|\bar{X}|]
\leq |2p_{ij}-1| + \E[|\bar{X} - \E[\bar{X}]|]
\leq |2p_{ij}-1| + \sqrt{\mathrm{Var}(\bar{X})}
= |2p_{ij}-1| + O(W^{-1/2})$.

Under Gaussian noise: $p_{ij} = \Phi(\mu_{ij}/\sigma_{ij})$.
For small SNR: $|2\Phi(x)-1| \leq \sqrt{2/\pi}\,|x| < |x|$ for $|x| \leq 1$.

Hoeffding on the bounded $X_m \in \{-1,+1\}$:
$\Pr(|\bar{X}| \geq \tau_s)
\leq 2\exp(-W(\tau_s-|2p_{ij}-1|)^2/2)$ when $\tau_s > |2p_{ij}-1|$.

\textbf{Step 2: Per-entry regret from updating.}
Updating entry $(i,j)$ of the projection matrix changes the
tracking error as:
$\hat{P}_{ij}^{\mathrm{new}} = \hat{P}_{ij}^{\mathrm{old}} - \eta_s G_{ij}$.
The expected squared error change (before retraction):
\begin{align*}
  \E[(\hat{P}_{ij}^{\mathrm{new}} - P_{ij})^2]
  &= e_{ij}^2 - 2\eta_s\mu_{ij}e_{ij} + \eta_s^2(\mu_{ij}^2+\sigma_{ij}^2).
\end{align*}
The improvement $\Delta R_{ij} = 2\eta_s\mu_{ij}e_{ij}
- \eta_s^2(\mu_{ij}^2+\sigma_{ij}^2)$.

\textbf{Step 3: When skipping helps.}
Skipping entry $(i,j)$ is beneficial when $\Delta R_{ij} \leq 0$, i.e.,
$2|\mu_{ij}||e_{ij}| \leq \eta_s(\mu_{ij}^2+\sigma_{ij}^2)$.
This is the state-dependent condition stated in the proposition.
Summing over all skipped entries and $T/B$ periods gives the bound.
\end{proof}

\begin{remark}[When selective consolidation can fail]
\label{rem:selective-fail}
If the threshold $\tau_s$ is set too aggressively (too high), signal
coordinates may be misclassified as transient, causing the learner to
miss genuine subspace drift. The synthetic experiment E4 showed this
failure mode with the previous covariance-threshold criterion, which
was too sensitive to the noise distribution.
The gradient-sign criterion is more robust because it directly measures
the consistency of the optimization direction.
\end{remark}


\bibliographystyle{plain}
% \bibliography{refs}

\end{document}
